% anexo_a.tex

\chapter{Anexos}
\label{anexos}

\section{Notación del modelo}
\label{anexo:notacion}

La Tabla~\ref{tab:notacion} reúne los símbolos empleados en el Capítulo~\ref{c2}
y en el Capítulo~\ref{c3}, con el significado que cada uno tiene en el modelo de
\citeA{jalan_incentive-aware_2024} y en su extensión dual de
\citeA{venegas_reyes_financial_2025}.

\begin{table}[htbp]
\centering
\caption[Notación del modelo de redes duales]{Notación empleada en el documento para el modelo de redes duales y su calibración empírica.}
\label{tab:notacion}
\small
\begin{tabular}{ll}
\hline
Símbolo & Significado \\
\hline
$W$ & Matriz simétrica de tamaños de contrato. $W_{ij}$ es el contrato vigente entre $i$ y $j$ \\
$P$ & Matriz antisimétrica de precios o pagos laterales entre pares de agentes \\
$M$ & Matriz de creencias de retorno esperado. $\mu_i = Me_i$ es el vector de creencias de $i$ \\
$\Sigma_i$ & Matriz de covarianza de retornos percibida por el agente $i$ \\
$\gamma_i$ & Coeficiente de aversión al riesgo del agente $i$ \\
$\Gamma$ & Matriz diagonal que agrupa los $\gamma_i$ de todos los agentes \\
$\Psi_i$ & Matriz que codifica qué contrapartes tiene permitidas el agente $i$ \\
$Q_i$ & Matriz de sensibilidad del agente $i$, definida en la Ecuación~\eqref{eq:optimo} \\
$F$ & Conjunto de pares $(i,j)$ con $i<j$ donde el precio $P_{ij}$ puede ser no nulo \\
$Z_F$ & Matriz de coeficientes del sistema de precios de la Ecuación~\eqref{eq:teorema1} \\
$\eta$ & Factor de amortiguación del algoritmo de negociación bilateral, Ecuación~\eqref{eq:eta} \\
$\lambda_A$ & Parámetro de fricción de acceso a la red Alternativa (modelo dual) \\
$\tau$ & Parámetro de transparencia o fricción del modelo dual \\
$m$ & Número de observaciones históricas por agente para estimar $\Sigma_i$ \\
$n$ & Número de agentes de la red. En el caso de aplicación, $n=4$ \\
\hline
\end{tabular}
\end{table}

\section{Fuentes de datos consultadas}
\label{anexo:fuentes}

La Tabla~\ref{tab:fuentes} detalla las fuentes revisadas durante el
levantamiento descrito en la Sección~\ref{c3:datos}, el régimen de acceso de
cada una y el insumo que aporta a la calibración del modelo.

\begin{table}[htbp]
\centering
\caption[Fuentes de datos consultadas para la calibración empírica]{Fuentes de datos consultadas, régimen de acceso e insumo que aporta cada una a la calibración del caso CORFO--Intermediario--MiPyME--Banco.}
\label{tab:fuentes}
\small
\begin{tabular}{p{0.26\textwidth}p{0.12\textwidth}p{0.50\textwidth}}
\hline
Fuente & Acceso & Insumo que aporta \\
\hline
Cuentas Públicas Participativas de CORFO, colección 2014--2024 & Público &
Monto de préstamos de CORFO a los IFNB ($W_{12}$) para 2019, 2020 y 2024, y
cifra de colocaciones del programa antecesor Microcrédito para 2016 \\
DIPRES, Informe de Monitoreo del programa 2020--2021 & Público &
Tasa de interés mayorista CORFO$\to$IFNB y spread de la tasa a cliente final
respecto del benchmark de mercado ACHEF/EFA \\
ELE-7, microdatos \texttt{ele7-full.csv} & Público &
Tasa de interés anualizada del crédito de mayor monto aprobado en 2022, por
tamaño de empresa, empleada como aproximación de $M_{34}$ \\
ELE-7, informe oficial de resultados & Público &
Tasa de rechazo condicional a la solicitud de crédito por tamaño de empresa,
empleada como ancla del supuesto sobre $\lambda_A$ \\
Reglamento operativo del programa Crédito Corfo Mipyme, Texto Refundido, punto 19 & No público &
Tasa de interés y saldo de capital vigente por operación y por mes, que
resolvería directamente el vacío de $W_{23}$ \\
\hline
\end{tabular}
\end{table}
